%%
%% This is file `ex_apresentacao.tex',
%% generated with the docstrip utility.
%%
%% The original source files were:
%%
%% UnBeamer.dtx  (with options: `presentation')
%% 
%% Arquivo gerado automaticamente.
%% 
%% Copyright (C) 2012 by Guilherme N. Ramos (gnramos@unb.br)
%% 
%% Este arquivo pode ser distribuído e/ou modificado conforme:
%%     1. LaTeX Project Public License e/ou
%%     2. GNU Public License.

%% Exemplo de utilização do tema UnB.
\documentclass{UnBeamer}%

\usepackage[brazilian]{babel}%
\usepackage[utf8]{inputenc}%
\usepackage{lmodern}%

%% ------------ Autor, Título, etc. ------------ %%
\title[Beamer + UnB]{Tema/Classe para apresentações utilizando Beamer}%
\subtitle{Exemplo Funcional}%
\author[gnramos]{Guilherme N. Ramos\\%
                 \texttt{gnramos@unb.br}}%
\institute[CIC]{Departamento de Ciência da Computação\\%
                Universidade de Brasília}%
\date[2012]{}%

%% Formata um comando.
\makeatletter%
\newcommand{\comm}[2][]{%
    \texttt{{\textbackslash}#2}%
    \if#1\@empty%
    \else\{#1\}\fi%
}%
\makeatother

%% ------------ Documento ------------ %%
\begin{document}%
  \frame{\titlepage}%

  \section{Tema UnB}%
  \begin{frame}%
    \begin{itemize}%
      \item Simples%
      \item Definição automática:%
        \begin{itemize}%
          \item título: nome da [sub]seção%
          \item rodapé: versão curta do título - nome da seção%
        \end{itemize}%
      \item Rodapé%
        \begin{description}%
          \item[\comm{unbeamerleftfooter}] instituto%
          \item[\comm{unbeamercentralfooter}] título resumido - nome da seção%
          \item[\comm{unbeamerrightfooter}] número de páginas/total de páginas%
        \end{description}%
      \item Definição de cores %
      \item Definição de listas%
    \end{itemize}%
  \end{frame}%

  \section{Listas \& Cores}%
  \begin{frame}%
    \begin{columns}[T]%
      \column{.45\textwidth}%
        Itens:\\%
        \begin{itemize}%
          \item {\color{redUnB}redUnB}%
          \begin{itemize}%
            \item {\color{greenUnB}greenUnB}%
            \begin{itemize}%
              \item {\color{blueUnB}blueUnB}%
            \end{itemize}%
          \end{itemize}%
        \end{itemize}%
      \column{.45\textwidth}%
        Enumerações:%
        \begin{enumerate}%
          \item \colorbox{redUnB}{\color{white}{redUnB}}%
          \begin{enumerate}%
            \item \colorbox{greenUnB}{\color{white}{greenUnB}}%
            \begin{enumerate}%
              \item \colorbox{blueUnB}{\color{white}{blueUnB}}%
            \end{enumerate}%
          \end{enumerate}%
        \end{enumerate}%
    \end{columns}%
    \vfill%
    Descrições:%
    \begin{description}%
      \item[Beamer] é uma classe.%
      \item[UnBeamer] é outra classe que redefine algumas características%
      para apresentações em Beamer%
    \end{description}%
  \end{frame}%

  \section{UnBeamer}%
  \subsection{Blocos}%
  \begin{frame}%
    \framesubtitle{(para a classe UnBeamer)}%
    \begin{alertblock}[.5\textwidth]{alertblock}%
      E agora os blocos%
    \end{alertblock}%
    \vfill\pause%
    \begin{exampleblock}[.75\textwidth]{exampleblock}%
      destacam a informação%
    \end{exampleblock}%
    \vfill\pause%
    \begin{block}{block}%
      com largura [opcional] variável.%
    \end{block}%
  \end{frame}%

  \renewcommand{\unbeamerleftfooter}{\colorbox{redUnB}{\color{white}{\secname}}}%

  \subsection{Comandos}%
  \begin{frame}%
    Mudança da parte central do rodapé (que continua atualizado com o
    nome da seção).%
    \subitem[$\rightarrow$]{%
        \comm[\protect{\comm{unbeamerleftfooter}}]{renewcommand}%
        \footnote[frame]{Exemplo de uso do comando \comm{subitem}.}}%
    \vfill%
    \begin{columns}%
      \column{.35\textwidth}%
        Inclusão de um pouco de texto...\\\vfill%
        ... e de uma imagem \strikeText{para ocupar espaço}\footnote[frame]{Exemplo%
        de uso do comando \comm{strikeText}.}.%
      \column{.45\textwidth}%
        \includegraphics[width=\textwidth]{UnBMarca}%
    \end{columns}%
    \vfill%
    E no slide seguinte, o resultado do comando \comm{imageFrame}.%
  \end{frame}%

  \imageFrame{UnBMarca}%
\end{document}%
\endinput
%%
%% End of file `ex_apresentacao.tex'.
